\documentclass[11pt,a4paper]{article}
\usepackage[margin=1in]{geometry}
\usepackage[T1]{fontenc}
\usepackage{lmodern}
\usepackage{amsmath,amssymb,amsthm}
\usepackage{graphicx}
\usepackage{microtype}
\usepackage{xcolor}
\usepackage[colorlinks=true,linkcolor=blue!55!black,citecolor=blue!55!black,urlcolor=blue!55!black]{hyperref}

\newtheorem{theorem}{Theorem}
\newtheorem{corollary}[theorem]{Corollary}
\newtheorem{remark}[theorem]{Remark}
\newcommand{\R}{\mathbb R}
\newcommand{\cL}{\mathcal L}

\title{Question 6 beyond Hilbert spaces:\\a Hilbert-separating factor criterion}
\author{Autonomous mathematical discovery packet\\
Run \texttt{fa\_banach\_001}, lane 0, \texttt{agent\_lane\_00}}
\date{11 August 2026}

\begin{document}
\maketitle

\begin{center}
\fbox{\parbox{0.91\textwidth}{\textbf{Status.}
Substantial partial result, likely valid, pending human review.  The arbitrary
nonseparable case remains open.}}
\end{center}

\section{Source question}

Thieu and Yen~\cite{TY} ask the following.  Let $E$ be a real normed space
and let $F\subset E$ be an infinite-dimensional linear subspace.  Do there
exist a continuous map $\psi:E\to E^*$ and a bounded linear map
$\varphi:E\to E^*$ such that
\begin{align*}
 &\text{for every }y\in E\text{ some }x\in F\text{ has }\psi(y)(x)\ne0,
 \tag{a3}\\
 &\text{for every }x\in F\setminus\{0\}\text{ some }y\in E\text{ has }
 \varphi(y)(x)\ne0,
 \tag{b3}
\end{align*}
and
\begin{equation}\label{eq:disp}
 \sup_{y\in E}\|\varphi(y)-\psi(y)\|_{E^*}<\infty?
\end{equation}
They prove this when $E$ is Hilbert and state that the non-Hilbert case
remains open.

\begin{figure}[ht]
\centering
\includegraphics[width=0.94\textwidth]{figures/open_problem_crop.png}
\caption{Question 6 in arXiv:2502.00746v1, page 20.}
\end{figure}

We use their following theorem~\cite[Theorem 4.5]{TY}: every
infinite-dimensional normed space $M$ admits a continuous map
$h:M\to M$ such that
\begin{equation}\label{eq:source}
 h(m)\ne0\quad(m\in M),\qquad
 \sup_{m\in M}\|h(m)-m\|\le1.
\end{equation}

\section{Proof intuition}

The Hilbert structure is needed only to manufacture a positive bilinear form,
not on the ambient space itself.  If a bounded operator sends $E$ into a
Hilbert space and is injective on $F$, then
$\varphi=T^*T$ has two key properties.  It separates every nonzero vector of
$F$, and no nonzero element of its range can vanish identically on $F$.
The range $M=\varphi(E)$ is therefore a normed space whose nonzero vectors are
already good for condition (a3).  Applying \eqref{eq:source} inside $M$ and
composing with $\varphi$ removes the sole zero without changing $\varphi$ by
more than one.

\section{The Hilbert-factor criterion}

\begin{theorem}\label{thm:criterion}
Let $E$ be a real normed space and $F\subset E$ an infinite-dimensional
linear subspace.  Suppose there are a real Hilbert space $H$ and a bounded
linear operator $T_0:E\to H$ whose restriction to $F$ is injective.  Then
Question 6 has an affirmative answer.  In fact, $\varphi$ and $\psi$ can be
chosen so that the supremum in \eqref{eq:disp} is at most $1$.
\end{theorem}

\begin{proof}
Let $G=\overline{T_0(F)}\subset H$, let $P_G:H\to G$ be the orthogonal
projection, and put $T=P_GT_0:E\to G$.  For $x\in F$ one has
$T_0x\in G$, hence $Tx=T_0x$.  Thus $T|_F$ is injective and $T(F)$ is dense
in $G$.

Define
\[
 \varphi=T^*T:E\longrightarrow E^*.
\]
For $x,y\in E$,
\begin{equation}\label{eq:gram}
 \varphi(y)(x)=\langle Ty,Tx\rangle_G.
\end{equation}
If $x\in F\setminus\{0\}$, then
\[
 \varphi(x)(x)=\|Tx\|_G^2>0.
\]
This proves (b3), with $y=x$.  It also shows that $\varphi|_F$ is injective,
so the normed subspace
\[
 M:=\varphi(E)\subset E^*
\]
is infinite-dimensional.

We next record the decisive range property:
\begin{equation}\label{eq:restriction}
 m\in M\setminus\{0\}\quad\Longrightarrow\quad m|_F\ne0.
\end{equation}
Indeed, write $m=\varphi(y)$.  If $m(x)=0$ for every $x\in F$, then
\eqref{eq:gram} says that $Ty$ is orthogonal to $T(F)$.  Since $Ty\in G$ and
$T(F)$ is dense in $G$, this forces $Ty=0$, whence
$m=T^*Ty=0$.  This proves \eqref{eq:restriction}.

Apply the source theorem \eqref{eq:source} to the infinite-dimensional normed
space $M$ (no closedness of $M$ is required) and obtain a continuous
nonvanishing $h:M\to M$ with $\|h(m)-m\|\le1$.  Set
\[
 \psi=h\circ\varphi:E\longrightarrow M\subset E^*.
\]
For every $y\in E$, $\psi(y)$ is a nonzero member of $M$, so
\eqref{eq:restriction} supplies an $x\in F$ with $\psi(y)(x)\ne0$.  Thus
(a3) holds.  Finally,
\[
 \|\psi(y)-\varphi(y)\|_{E^*}
 =\|h(\varphi(y))-\varphi(y)\|_M\le1,
\]
which proves \eqref{eq:disp} with the asserted bound.
\end{proof}

\section{Countably total and separable subspaces}

A family $(f_n)\subset F^*$ is \emph{total} if
$f_n(x)=0$ for every $n$ implies $x=0$.

\begin{corollary}\label{cor:total}
Question 6 has an affirmative answer whenever $F^*$ contains a countable
total family.  In particular, it has an affirmative answer whenever $F$ is
separable.  The subspace $F$ need not be closed or complemented, and $E$ need
not be complete or Hilbert.
\end{corollary}

\begin{proof}
Let $(f_n)_{n\ge1}\subset F^*$ be total.  By Hahn--Banach, extend each $f_n$
to $\widetilde f_n\in E^*$ without changing its norm.  Put
\[
 a_n=\frac{2^{-n}}{1+\|\widetilde f_n\|},\qquad
 T_0y=(a_n\widetilde f_n(y))_{n\ge1}\in\ell_2.
\]
Then
\[
 \|T_0y\|_2^2\le \|y\|^2\sum_{n\ge1}4^{-n},
\]
so $T_0:E\to\ell_2$ is bounded.  Its restriction to $F$ is injective because
all $a_n$ are nonzero and $(f_n)$ is total.  Theorem~\ref{thm:criterion}
applies.

If $F$ is separable, choose a dense sequence $(u_n)$ in its unit sphere and,
by Hahn--Banach, choose $f_n\in F^*$ with $\|f_n\|=f_n(u_n)=1$.  For nonzero
$x\in F$, some $u_n$ lies within $1/2$ of $x/\|x\|$, and then
\[
 f_n(x/\|x\|)>f_n(u_n)-\tfrac12=\tfrac12.
\]
Thus $(f_n)$ is total.
\end{proof}

\begin{remark}
Theorem~\ref{thm:criterion} also recovers the source's Hilbert case, including
nonclosed $F$: take $H=E$ and $T_0$ to be the orthogonal projection onto
$\overline F$, which restricts to the identity on $F$.
\end{remark}

\section{Scope, upgrade obstruction, and novelty}

The theorem covers all separable $F$ and also nonseparable spaces possessing
a countable total dual family.  For a completely arbitrary nonseparable
$F$, an uncountable total family does not automatically yield a bounded map
into $\ell_2(I)$: a summable positive weight family has countable support,
while unweighted evaluations need not be square summable.  More generally,
(b3) already requires a bounded bilinear form that is right-nondegenerate on
$F$.  No universal construction or explicit counterexample was found, so the
full arbitrary-$F$ question remains open.

A bounded search through 11 August 2026 covered the run indexes; exact arXiv
id, title, question, and authors; follow-up and citation queries; non-Hilbert
maximum displacement; Hilbert-separating operators; countable total dual
families; and nondegenerate bounded bilinear forms.  No later answer or this
criterion was found.

\begin{thebibliography}{9}
\bibitem{TY}
N. N. Thieu and N. D. Yen,
\emph{The Hartman--Stampacchia Theorem and the Maximum Displacements of
Nonvanishing Continuous Vector-Valued Functions}, arXiv:2502.00746v1 (2025),
especially Question 6 and Theorems 4.5--4.6.
\end{thebibliography}

\end{document}

